A tuned LLM system is usually reported with its score on the same benchmark items that were used to search for its prompt or program. We ask how far such a score is from performance on fresh items, what drives the gap, whether it changes conclusions, and how to report more reliably. In a synthetic world with exact truth (203{,}400 search replicates) and on a real model and benchmark (Qwen2.5-1.5B-Instruct on MMLU-Pro, 300 candidate configurations replayed from cached outputs), the naive report is optimistic by up to 14.6 points (synthetic, 50 items, 200 candidates, no true differences) and 5.7 points on MMLU-Pro (100 search items, 300 candidates), roughly doubling the apparent gain over the baseline; the gap shrinks like $1/\sqrt n$, grows with the number of candidates and few-shot demonstrations drawn from the search set, and follows an extreme-value formula whose parameters we fit to each setting (a check of its form, not a held-out prediction). Naive reports declare a tuned system better than an untuned one in 97.6\% of synthetic replicates with no true gain; on MMLU-Pro such false-claim rates are not identifiable (the baseline is almost always truly worse), so decision errors rest on the synthetic world and a paired A-versus-B design on real outcomes, where reports rank systems by search effort. Hold-out and conditional/hybrid selection-aware intervals restore nominal coverage but cost items, width and power; empirical-Bayes disclosure fails. Results rest on one small model, template and ensemble configurations replayed from cached outputs on one fixed 2000-item pool (between-pool uncertainty is not estimated), and no live tool run.
